
@@M0@@ de @@M1@@ dans @@M2@@ tel que @@M3@@.

1.3. On définit pour tout @@M4@@ une filtration sur @@M5@@ de la façon suivante : une série @@M6@@ sera dite de filtration @@M7@@ si @@M8@@ pour @@M9@@. On note @@M10@@ le sous-groupe ainsi obtenu ; c'est un foncteur sur @@M11@@, et nous considérerons désormais @@M12@@ comme foncteur à valeurs dans la catégorie @@M13@@ des groupes abéliens filtrés grâce à cette filtration(avec @@M14@@).

Soit @@M15@@ un foncteur sur @@M16@@ à valeurs dans @@M17@@~; nous supposerons que @@M18@@ satisfait les deux conditions suivantes :
@@M19@@

En fait, nous n'aurons besoin de ces conditions que lorsque @@M20@@ est un anneau de polynômes, et @@M21@@ un produit fini de groupes symétriques opérant par permutation des indéterminées.

\textbf{Exemples}~: 1)~@@M22@@. 2)~@@M23@@~; ce foncteur sera noté @@M24@@ (@@M25@@, initiale de « grossière »). 3)~@@M26@@, où @@M27@@ désigne le groupe (pour l'addition) des séries formelles telles que @@M28@@, @@M29@@ pour @@M30@@, muni de la filtration analogue à celle de @@M31@@.

1.4. Pour tout @@M32@@, nous noterons @@M33@@ l'endomorphisme de @@M34@@
